% Section 4.4 -- storage rule on simulated scenarios (condensed 28 Oct 2026).
\subsection{Storage Rule on Simulated Scenarios}
\label{sec:results-storage}

The node has not run in real grain, so the rule was tested on simulated streams for basmati paddy stored near Patiala.
Air temperature and humidity follow the monthly means of the IMD station at Patiala for 1991--2020~\cite{imd2022} (08:30 and
17:30 IST for humidity); shed air is assumed equal to outdoor air, grain temperature follows it with a 2-day lag and water
activity drifts slowly towards the shed's daily mean humidity. These assumptions are listed in the repository. The
\kn{stor_scenarios} scenarios cover dry grain held from November to February, a monsoon hold, a lot probed straight after
harvest and sun-dried on day 8, the monsoon with Wi-Fi and power outages, a probe hanging in shed air, and grain at
$a_w\approx0.65$; each ran with \kn{stor_seeds} noise seeds and with sensor biases of $\pm2$\,\%RH. STORAGE-DEMO-0.1 was compared
with three simpler rules (latest reading only, shed-air humidity, temperature alone), each scored against the category that the
rule's definition gives for the noise-free grain at the same hour. This measures how far noise, outages and shortcuts move a
rule from its definition; it does not show that the definition predicts spoilage.

Without bias, STORAGE-DEMO-0.1 agreed with the intended category in \kn{stor_s3_demo_agree}--\kn{stor_s2_demo_agree}\,\% of known
hours in the four in-grain scenarios and changed category at most \kn{stor_demo_flips_max} times a week. The latest-reading rule
agreed about as often but changed \kn{stor_instant_flips_min}--\kn{stor_instant_flips_max} times a week, and shed-air humidity was a poor
stand-in for the grain (agreement \kn{stor_shed_agree_min}--\kn{stor_shed_agree_max}\,\%). Temperature alone gave no HIGH hour in the winter
hold, when moisture rather than heat was the problem. The costs are delay and unknowns: the 24-hour median reported HIGH about
\kn{stor_s2_demo_highdelay}\,h after the intended category, the first day after insertion is UNKNOWN, and with the outages the
rule was UNKNOWN for \kn{stor_s4_demo_unknown}\,\% of the hours against \kn{stor_s4_instant_unknown}\,\% for the latest-reading
rule. A sensor error of $\pm2$\,\%RH mattered about as much as the choice of rule: with $+2$\,\%RH up to
\kn{stor_demo_falsehigh_biasplus_max}\,\% of hours were a false HIGH and with $-2$\,\%RH up to \kn{stor_demo_missedhigh_biasminus_max}\,\%
a missed HIGH, so a salt calibration~\cite{greenspan1977} is required before real readings are interpreted.

\begin{figure}[t]
  \centering
  \includegraphics[width=0.95\textwidth]{fig_storage_scenarios.pdf}
  \caption{Simulated storage scenarios (seed 0): probe and shed humidity with the 65 and 70\,\% limits, and the category per
  hour for the intended definition, STORAGE-DEMO-0.1 and three simpler rules. (a) Monsoon hold; (b) fresh lot, sun-dried on day
  8; (c) monsoon with outages (grey); (d) probe in shed air. Simulated data; categories describe conditions, not spoilage.}
  \label{fig:storage}
\end{figure}
